\subsection*{الف}
هیستوگرام تصویر $8\times4$ (در مجموع $N=32$ پیکسل، سطوح خاکستری $0..7$):
\begin{flalign*}
    & n_0=6,\; n_1=4,\; n_2=4,\; n_3=2,\; n_4=4,\; n_5=3,\; n_6=5,\; n_7=4 &&
\end{flalign*}
\begin{flalign*}
    p &= \left\{\tfrac{6}{32},\tfrac{4}{32},\tfrac{4}{32},\tfrac{2}{32},\tfrac{4}{32},\tfrac{3}{32},\tfrac{5}{32},\tfrac{4}{32}\right\} \\
    H &= -\sum_{i} p_i \log_2 p_i = 2.941 \ \ \text{\lr{bits/pixel}} &&
\end{flalign*}

\subsection*{ب}
درخت
\lr{Huffman}
با ادغام پیاپی کم‌احتمال‌ترین‌ها، به کدهای زیر می‌رسد:
\begin{flalign*}
    & 0:00,\quad 6:111,\quad 1:010,\quad 2:011,\quad 4:100,\quad 7:101,\quad 5:1101,\quad 3:1100 &&
\end{flalign*}
طول متوسط، کارایی و افزونگی کدگذاری:
\begin{flalign*}
    L_{avg} &= \sum_i p_i\, l_i = \tfrac{1}{32}\big(6\cdot2+5\cdot3+4\cdot3+4\cdot3+4\cdot3+4\cdot3+3\cdot4+2\cdot4\big) \\
            &= \tfrac{95}{32} = 2.969 \ \ \text{\lr{bits/pixel}} \\
    \eta &= \frac{H}{L_{avg}} = \frac{2.941}{2.969} = 0.991 \\
    R &= L_{avg} - H = 0.027 \ \ \text{\lr{bits/pixel}} &&
\end{flalign*}

\subsection*{ج}
بازده $99.1\%$ است؛ نسبت به کدگذاری ثابت $3$ بیتی صرفه‌جویی ناچیزی به‌دست می‌آید:
\begin{flalign*}
    C = \frac{3}{2.969} = 1.01 &&
\end{flalign*}
علت ضعف کارایی آن است که توزیع احتمال سطوح خاکستری تقریباً یکنواخت است؛ کدگذاری
\lr{Huffman}
تنها افزونگی کدگذاری
(\lr{coding redundancy})
را حذف می‌کند و وقتی مؤثر است که هیستوگرام بسیار غیریکنواخت و چوله باشد—یعنی چند سطح با احتمال بسیار بالا (نزدیک به توان‌های منفی $2$) و بقیه با احتمال بسیار کم. تنها در این حالت طول‌های متغیر کوتاه به سطوح پرتکرار اختصاص یافته و $L_{avg}\to H$ کاهش می‌یابد.

\subsection*{د}
افزونگی واقعی این تصویر، افزونگی مکانی
(\lr{interpixel redundancy})
است نه افزونگی کدگذاری؛ پس پیش از
\lr{Huffman}
از کدگذاری تفاضلی
\lr{DPCM}
استفاده می‌کنیم: در هر سطر پیکسل اول حفظ و بقیه با $d=f_i-f_{i-1}$ جایگزین می‌شوند. تصویر تفاضلی:
\begin{flalign*}
    & 0,0,2,0,5,0,0,0 \\
    & 0,0,2,-2,3,3,0,-2 \\
    & 0,2,-1,2,3,0,-1,-1 \\
    & 1,0,0,5,-1,0,-1,0 &&
\end{flalign*}
هیستوگرام تفاضل بسیار چوله است ($0$ پانزده بار تکرار شده) و آنتروپی کاهش می‌یابد:
\begin{flalign*}
    & p(0)=\tfrac{15}{32},\, p(-1)=\tfrac{5}{32},\, p(2)=\tfrac{4}{32},\, p(3)=\tfrac{3}{32},\, p(-2)=p(5)=\tfrac{2}{32},\, p(1)=\tfrac{1}{32} \\
    H_d &= 2.282 \ \ \text{\lr{bits/pixel}} &&
\end{flalign*}
کد
\lr{Huffman}
روی تفاضل‌ها:
\begin{flalign*}
    & 0:0,\quad -1:110,\quad 2:101,\quad 3:100,\quad 5:1110,\quad -2:11111,\quad 1:11110 &&
\end{flalign*}
\begin{flalign*}
    L_{avg}' &= \tfrac{1}{32}\big(15\cdot1+5\cdot3+4\cdot3+3\cdot3+2\cdot4+2\cdot5+1\cdot5\big) = \tfrac{74}{32}=2.313 \\
    \eta' &= \frac{2.282}{2.313}=0.987,\qquad C'=\frac{3}{2.313}=1.30 &&
\end{flalign*}

\subsection*{ه}
روش بدون اتلاف دیگر، کدگذاری طول اجرا
(\lr{Run-Length Encoding})
روی پیمایش سطری است، که افزونگی مکانی ناشی از تکرار پیکسل‌های همسایه را هدف می‌گیرد. زوج‌های $(\text{\lr{gray}},\text{\lr{run}})$:
\begin{flalign*}
    & (0,2)(2,2)(7,4)(0,2)(2,1)(0,1)(3,1)(6,2)(4,1)(0,1) \\
    & (2,1)(1,1)(3,1)(6,2)(5,1)(4,1)(1,3)(6,1)(5,2)(4,2) &&
\end{flalign*}
\begin{flalign*}
    & 32 \ \text{پیکسل} \longrightarrow 20 \ \text{زوج} &&
\end{flalign*}
نوع افزونگی هدف‌گرفته‌شده افزونگی مکانی
(\lr{interpixel / spatial redundancy})
است؛ با ترکیب خروجی
\lr{RLE}
و کدگذاری
\lr{Huffman}
روی سطوح و طول اجراها، افزونگی کدگذاری باقی‌مانده نیز حذف می‌شود.